% Generated draft: requires mathematical and source review.
\documentclass[11pt]{article}
\usepackage[margin=25mm]{geometry}
\usepackage{fontspec}
\setmainfont{Latin Modern Roman}
\usepackage{amsmath,amssymb,mathtools,unicode-math}
\setmathfont{Latin Modern Math}
\usepackage{enumitem,needspace,xurl,hyperref}
\setlength{\parindent}{0pt}
\setlength{\parskip}{5pt}
\setlength{\emergencystretch}{3em}
\newcommand{\N}{\mathbb{N}}
\newcommand{\Z}{\mathbb{Z}}
\newcommand{\Q}{\mathbb{Q}}
\newcommand{\R}{\mathbb{R}}
\newcommand{\C}{\mathbb{C}}
\newcommand{\F}{\mathbb{F}}
\newcommand{\A}{\mathbb{A}}
\newcommand{\PP}{\mathbb{P}}
\newcommand{\E}{\mathbb{E}}
\newcommand{\eps}{\varepsilon}
\newcommand{\dd}{\,\mathrm{d}}
\begin{document}
\section*{UnsolvedMath ID: 800153}
\textbf{Source problem code:} catalogue entry
\subsection{Definitions and mathematical statement}
Let \(\log\) denote the natural logarithm. The Riemann hypothesis asserts that every nontrivial zero of the Riemann zeta function \(\zeta(s)\) has real part \(1/2\).

The source denotes a zeta-moment quantity by \(I_k(T)\), but the supplied record does not define it. In particular, its integration interval and normalization are not supplied. The statements below therefore retain this notation as an unresolved source-defined quantity; they require its exact definition before constituting fully specified mathematical statements.

\textbf{Requested explicit bound.} Assuming the Riemann hypothesis, determine, for every real \(k>0\), an explicit constant \(C_k>0\) such that there exists a real threshold \(T_0(k)>1\) for which
\[
\forall T\in\R,\quad T\ge T_0(k)\ \Longrightarrow\ I_k(T)\le C_k T(\log T)^{k^2}.
\]
This is the requested explicit version of Harper's upper bound, referenced in the source as [Ha]. The source asks for an explicit \(C_k\), but does not expressly require an explicit threshold \(T_0(k)\).

\textbf{Related conjecture quoted by the source.} The source attributes to Soundararajan [S2] the assertion
\[
\forall k\in\R,\ k\ge2,\quad\forall T\in\R,\ T\ge10,\qquad
I_k(T)\le T(\log T)^{k^2}.
\]
This is a separately quoted conjecture, not an established bound. The supplied excerpt does not explicitly specify whether this related conjecture is conditional on the Riemann hypothesis.
\subsection{Short English statement}
Assuming the Riemann hypothesis, find an explicit constant for each positive real moment order in Harper's upper bound for sufficiently large \(T\). The source also quotes a conjectural constant-one bound of Soundararajan for \(k\ge2\) and \(T\ge10\). The definition of \(I_k(T)\) is missing from the supplied record and must be recovered before mathematical review.
\subsection{Sources}
\begin{itemize}
\item Local UnsolvedMath registry: permanent ID \(800153\), source problem code \texttt{catalogue entry}, titled ``Moments of L-functions --- Moments of the Riemann zeta function --- Problem 1.'' The supplied registry labels this problem open; that label is not an independently verified assessment.
\item Hamieh, Kadiri, Ng, \emph{Moments of L-functions Problem List} (2024), Problem 1: \url{https://arxiv.org/abs/2405.17800}. The supplied provenance records access on 2026-10-16 and extraction by \texttt{arxiv\_latex}; these are source-provided provenance statements, not claims of access for this draft.
\item The problem text credits Nathan Ng, Lethbridge. The registry lists Hamieh, Kadiri, Ng as proposers and 2024 as the proposed year.
\item Harper [Ha]: the upper bound to be made explicit. The supplied record does not include the bibliographic entry corresponding to [Ha].
\item Soundararajan [S2]: the related conjecture quoted in the statement. The supplied record does not include the bibliographic entry corresponding to [S2].
\end{itemize}
\end{document}
